% ======================================================
%  Rozdział — Mieszane Wyzwania
%
%  GENERATED by tools/convert_deck.py from areas/5-mieszane-wyzwania.tex.
%  Edit the deck source and re-run, or convert this chapter to hand-maintained
%  and delete it from the script's AREA_META -- but do not edit both.
%
%  The `% -- NN --` comments are the deck's own global exercise numbers, kept
%  so a reader of either source can find the same exercise in the other. The
%  book itself numbers frames per chapter: within a chapter, frame N is
%  exercise N, because every frame carries exactly one exercise except the last.
% ======================================================

\chapter{Mieszane Wyzwania}

Prędkość, procenty, geometria, finanse i praca — zadania z życia, liczone w pamięci.

% -- 34 --
\begin{fr}
  \exercise[\CategoryBadge[MixColor!20]{Pr\k{e}dko\'s\'c}]{2}{2 min}{%
      Poci\k{a}g jedzie 120~km w 1{,}5~h.\par
      Jaka jest jego \'{s}rednia pr\k{e}dko\'s\'c?}
\end{fr}

% -- 35 --
\begin{fr}
  \ans{80~km/h}{$120 \div 1{,}5 = 80$}
  \exercise[\CategoryBadge[MixColor!20]{Procenty}]{2}{2 min}{%
      Pizza ma 8 kawa\l{}k\'{o}w, zjadasz 3.\par
      Jaki procent zostaje?}
\end{fr}

% -- 36 --
\begin{fr}
  \ans{62{,}5\%}{$\tfrac{5}{8} \times 100\% = 62{,}5\%$}
  \exercise[\CategoryBadge[MixColor!20]{Geometria}]{3}{3 min}{%
      Pok\'{o}j $5\,\text{m}\times 4\,\text{m}$.\par
      Kafelki $50\,\text{cm}\times 50\,\text{cm}$.\par
      Ile kafelk\'{o}w?}
\end{fr}

% -- 37 --
\begin{fr}
  \ans{80}{$10 \times 8 = 80$}
  \exercise[\CategoryBadge[MixColor!20]{Finanse}]{3}{3 min}{%
      1000~z\l{} przy 5\% rocznie.\par
      Procent sk\l{}adany. 2~lata.\par
      Ile masz?}
\end{fr}

% -- 38 --
\begin{fr}
  \ans{1102{,}50~z\l{}}{$1000 \times 1{,}05^2$}
  \exercise[\CategoryBadge[MixColor!20]{Praca}]{3}{5 min}{%
      3 robotnik\'{o}w robi prac\k{e} w 6 dni.\par
      Ile dni zajmie 2 robotnikom?}
\end{fr}

% The chapter's last frame carries the last answer alone.
\begin{fr}
  \ans{9 dni}{$3 \times 6 = 18$ robotnikodni $\div$ $2 = 9$.}
\end{fr}
